\section{Baseline Methods}\label{sec:method:baselines}

The baselines isolate the contribution of the dual-branch layout in \cref{sec:method:arch} and of the mode term $\lambda_{\mathrm{md}}\Loss_{\mathrm{md}}$ in \cref{eq:method:loss_total} relative to a single-stack encoder and to classical clustering without a sequence model.
Final configurations and the full list of comparisons are deferred to \cref{ch:experiments}; this section fixes only the design intent of each.

\begin{description}
  \item[Single-stack transformer.] A flat encoder of $N_{\mathrm{ss}}{=}8$ transformer-encoder layers with the same width, heads, and \gls{ffn} expansion as the trunk in \cref{sec:method:arch}, terminated by one linear head $\R^{d_{\mathrm{model}}}\!\to\!\R^{p}$ and trained on $\Loss_{\mathrm{em}}$ alone. Tests whether a single representation optimised purely for window-level deinterleaving also serves the mode clustering.

  \item[\Gls{hdbscan} on fixed window features.] No transformer; \gls{hdbscan} \parencite{mcinnes2017hdbscan} clusters a deterministic vector built from the scaled \glspl{pdw} of a window (vectorisation specified in \cref{ch:experiments}). Separates the gain from learned sequence context from what density clustering already recovers on hand-specified summaries.

  \item[External references.] $k$-means on raw window vectors and on autoencoder embeddings, \gls{dec}-style objectives, SwAV \parencite{caron2020swav} assignment prediction, and a closed-catalogue supervised classifier as an upper bound on the modes seen during training. The supervised classifier is reported only as a reference ceiling; it cannot accommodate the held-out modulations of the open-set experiments and is not used as the fielded mode predictor.
\end{description}
